\documentclass[a4paper, serif, 10pt]{moderncv}

%% moderncv
% color and style
\moderncvcolor{black}
\moderncvstyle{banking}

%% page layout/margins
\usepackage[top=2.5cm, bottom=2.5cm, left=1.5cm, right=1.5cm]{geometry}

\nopagenumbers{}

%% Babel config for French language
% ref:
% - https://latex3.github.io/babel/guides/locale-french.html
\usepackage[french]{babel}

%% fontspec
\usepackage{fontspec}

\IfFontExistsTF{Linux Libertine}{
  \setmainfont[Ligatures={TeX, Common}, Numbers=OldStyle]{Linux Libertine}
}{}
\IfFontExistsTF{Linux Libertine O}{
  \setmainfont[Ligatures={TeX, Common}, Numbers=OldStyle]{Linux Libertine O}
}{}

%% CTeX
% CJK support, used to show CJK characters in the resume
%
% - heading=false: disable ctex heading
% - punct=kaiming: use kaiming punctuations styles for CJK
% - scheme=plain: use plain scheme, do not override \normalsize font size
% - space=auto: space settings for CJK characters
%
% CJK fonts are resolved per locale language, e.g. Noto Serif/Sans CJK JP
% for Japanese resumes, with fallback to the next region if not installed
%
% ref:
% - http://ctan.mirrorcatalogs.com/language/chinese/ctex/ctex.pdf
\usepackage[UTF8, heading=false, punct=kaiming, scheme=plain, space=auto]{ctex}

\IfFontExistsTF{Noto Serif CJK SC}{
  \setCJKmainfont{Noto Serif CJK SC}
}{}
\IfFontExistsTF{Noto Sans CJK SC}{
  \setCJKsansfont{Noto Sans CJK SC}
}{}

%% line spacing
\usepackage{setspace}
\setstretch{1.125}

%% URL styling - use normal text instead of monospace
\urlstyle{same}

%% Patch \httplink and \httpslink to handle full URLs
% The original moderncv macros always prepend http:// or https:// to URLs.
% This causes issues when the URL already has a protocol (e.g., https://example.com)
% resulting in malformed URLs like https://https://example.com.
% This patch checks if the URL already contains "://" and uses it directly if so.
% Note: We use \str_set:Nx (x-expansion) to expand macros like \@homepage first.
\ExplSyntaxOn
\str_new:N \g_yamlresume_url_str

\RenewDocumentCommand{\httpslink}{O{}m}{
  \str_set:Nx \g_yamlresume_url_str {#2}
  \str_if_in:NnTF \g_yamlresume_url_str {://}
    { \url{#2} }
    { \url{https://#2} }
}

\RenewDocumentCommand{\httplink}{O{}m}{
  \str_set:Nx \g_yamlresume_url_str {#2}
  \str_if_in:NnTF \g_yamlresume_url_str {://}
    { \url{#2} }
    { \url{http://#2} }
}
\ExplSyntaxOff

\name{Camille Morel}{}
\title{Ingénieur Données et Plateformes}
\phone[mobile]{+33 6 12 34 56 78}
\email{camille.morel@dataforge.fr}
\homepage{https://camillemorel-data.fr}

\address{24 rue des Lilas, Lyon, Auvergne-Rhône-Alpes, France, 69003}{}{}

\extrainfo{{\small \faLinkedin}\ \href{https://www.linkedin.com/in/camille-morel-data}{@camille-morel-data} {} {} {} • {} {} {} 
{\small \faGithub}\ \href{https://github.com/camillemorel-data}{@camillemorel-data}}

\begin{document}

\maketitle

\section{Informations de base}

\cvline{}{Ingénieur données avec plus de 5 ans d'expérience dans la conception et l'exploitation de plateformes data à grande échelle. J'accompagne les équipes métier et produit dans la valorisation de leurs données, de l'ingestion à la mise à disposition dans des outils décisionnels.
%
\begin{itemize}
\item Conception de pipelines ETL/ELT robustes et orchestrés avec Airflow et dbt
\item Modélisation d'entrepôts de données et de lacs de données sur Snowflake, BigQuery et AWS
\item Mise en place de pratiques DataOps : tests automatisés, CI/CD, observabilité et gouvernance
\item Collaboration étroite avec les scientifiques des données, analystes et équipes plateforme
\item Encadrement et mentorat de développeurs juniors
\end{itemize}}
\section{Formation}

\cventry{sept. 2016–juil. 2021}
        {Master, Informatique et Science des Données, Score : Mention Bien}
        {INSA Lyon}
        {\url{https://www.insa-lyon.fr}}
        {}
        {\begin{itemize}
\item Formation d'ingénieur orientée ingénierie des données, statistiques et systèmes distribués
\item Projet de fin d'études sur l'optimisation de requêtes dans un entrepôt de données cloud
\item Stage de 6 mois en équipe plateforme data chez un éditeur logiciel lyonnais
\end{itemize}
\textbf{Cours} : Bases de données avancées, Algorithmique distribuée, Apprentissage automatique, Entrepôts de données, Cloud computing, Statistiques appliquées, Architecture logicielle, Sécurité des données}
\section{Expérience professionnelle}

\cventry{mars 2023–Present}
        {Ingénieur Données Senior}
        {DataForge}
        {\url{https://www.dataforge.fr}}
        {}
        {\begin{itemize}
\item Conception et maintien de pipelines ELT orchestrés avec Airflow et dbt pour 12 sources hétérogènes
\item Modélisation d'un entrepôt Snowflake et mise en place de tests de qualité automatisés
\item Réduction de 35 \% du temps de traitement des flux critiques grâce à l'optimisation SQL et au parallélisme
\item Contribution à la gouvernance des données : catalogage, documentation et gestion des accès
\item Mentorat de deux ingénieurs données juniors et animation de revues de code hebdomadaires
\end{itemize}
\textbf{Mots-clés} : Airflow, dbt, Snowflake, DataOps, Gouvernance}

\cventry{sept. 2021–févr. 2023}
        {Ingénieur Données}
        {Groupe Voltaire}
        {\url{https://www.groupe-voltaire.fr}}
        {}
        {\begin{itemize}
\item Développement de flux d'ingestion temps réel avec Kafka et traitement batch avec Spark
\item Migration d'un entrepôt PostgreSQL vers BigQuery avec réduction des coûts de 20 \%
\item Industrialisation des déploiements via GitLab CI et conteneurs Docker
\item Participation à la définition du modèle de données et des indicateurs métier avec les analystes
\end{itemize}
\textbf{Mots-clés} : Python, SQL, Kafka, BigQuery, CI/CD}

\cventry{févr. 2021–août 2021}
        {Ingénieur Données Junior}
        {Lumière Analytics}
        {\url{https://www.lumiere-analytics.fr}}
        {}
        {\begin{itemize}
\item Création de jobs ETL Python pour alimenter un entrepôt PostgreSQL à partir de sources CRM et ERP
\item Automatisation de contrôles de cohérence et de rapports de qualité des données
\item Développement de tableaux de bord Power BI pour les équipes commerciales
\item Rédaction de la documentation technique et des procédures d'exploitation
\end{itemize}
\textbf{Mots-clés} : ETL, PostgreSQL, Docker, Power BI}
\section{Langues}

\cvline{Français}{Niveau natif ou bilingue \hfill \textbf{Mots-clés} : Langue maternelle}
\cvline{Anglais}{Niveau professionnel complet \hfill \textbf{Mots-clés} : TOEIC 950, Travail en équipe internationale}
\cvline{Espagnol}{Niveau de travail limité \hfill \textbf{Mots-clés} : Conversation professionnelle}
\section{Compétences}

\cvline{Ingénierie des données}{Expert \hfill \textbf{Mots-clés} : Python, SQL, dbt, Apache Airflow, Apache Spark, Kafka, Snowflake, BigQuery}
\cvline{Cloud \& DevOps}{Avancé \hfill \textbf{Mots-clés} : AWS, GCP, Docker, Kubernetes, Terraform, GitLab CI, GitHub Actions}
\cvline{Bases de données}{Avancé \hfill \textbf{Mots-clés} : PostgreSQL, Snowflake, BigQuery, MongoDB, Redis, Elasticsearch}
\cvline{Visualisation et BI}{Intermédiaire \hfill \textbf{Mots-clés} : Power BI, Tableau, Looker, Metabase}
\section{Récompenses}

\cventry{juin 2024}
        {DataForge}
        {Prix de l'impact data}
        {}
        {}
        {Récompense interne attribuée pour la refonte de la plateforme d'ingestion et l'amélioration significative de la fiabilité des données.}

\cventry{juil. 2021}
        {INSA Lyon}
        {Meilleur projet de fin d'études}
        {}
        {}
        {Projet sur l'optimisation de requêtes analytiques dans un entrepôt cloud, distingué pour sa rigueur technique et son impact pédagogique.}
\section{Certificats}

\cventry{avr. 2024}
        {Amazon Web Services}
        {AWS Certified Data Engineer - Associate}
        {\url{https://aws.amazon.com/certification/certified-data-engineer-associate/}}
        {}
        {}

\cventry{nov. 2023}
        {dbt Labs}
        {dbt Certified Developer}
        {\url{https://www.getdbt.com/certifications/}}
        {}
        {}

\cventry{sept. 2022}
        {Google Cloud}
        {Google Cloud Professional Data Engineer}
        {\url{https://cloud.google.com/certification/data-engineer}}
        {}
        {}
\section{Publications}

\cventry{févr. 2024}
        {Modernisation d'un entrepôt de données : retour d'expérience}
        {Revue Data \& Systèmes}
        {\url{https://www.revue-data-systemes.fr/articles/modernisation-entrepot}}
        {}
        {\begin{itemize}
\item Présentation d'une démarche de migration vers Snowflake menée en équipe
\item Retour sur les choix d'architecture, la qualité des données et la conduite du changement
\item Bonnes pratiques pour industrialiser les pipelines avec dbt et Airflow
\end{itemize}}
\section{Projets}

\cventry{janv. 2023–juin 2023}
        {Pipeline libre de valorisation des données publiques de la métropole lyonnaise}
        {OpenData Lyon}
        {\url{https://github.com/camillemorel-data/opendata-lyon}}
        {}
        {\begin{itemize}
\item Collecte automatisée de jeux de données ouvertes et stockage dans PostgreSQL
\item Transformations dbt et documentation du lignage des données
\item Tableau de bord public pour explorer les indicateurs de mobilité et de qualité de l'air
\end{itemize}
\textbf{Mots-clés} : Open Data, Airflow, dbt, PostgreSQL}

\cventry{mars 2022–déc. 2022}
        {Outil de surveillance de la qualité des données pour pipelines analytiques}
        {Moniteur Qualité des Données}
        {\url{https://github.com/camillemorel-data/data-quality-monitor}}
        {}
        {\begin{itemize}
\item Définition de règles de validation configurables et génération de rapports quotidiens
\item Intégration dans une chaîne CI/CD pour bloquer les déploiements en cas d'anomalie critique
\item Réduction du temps de détection des incidents de données de plusieurs heures à quelques minutes
\end{itemize}
\textbf{Mots-clés} : Qualité des données, Python, Great Expectations, Docker}
\section{Centres d'intérêt}

\cvline{Sport}{Course à pied, Escalade, Randonnée}
\cvline{Musique}{Guitare, Jazz}
\cvline{Communauté}{Rencontres professionnelles, Mentorat}
\section{Bénévolat}

\cventry{janv. 2022–déc. 2024}
        {Bénévole en ingénierie des données}
        {Data for Good Lyon}
        {\url{https://dataforgood.fr}}
        {}
        {\begin{itemize}
\item Participation à des projets data à impact social avec des associations locales
\item Conception de pipelines de collecte et de nettoyage de données ouvertes
\item Animation d'ateliers d'initiation à la donnée pour des bénévoles non techniques
\end{itemize}}
    

\cventry{sept. 2023–Present}
        {Marraine}
        {Elles bougent}
        {\url{https://www.ellesbougent.com}}
        {}
        {\begin{itemize}
\item Intervention dans des collèges et lycées pour promouvoir les carrières scientifiques et techniques auprès des jeunes filles
\item Organisation d'ateliers de découverte des métiers de la donnée et de l'ingénierie
\end{itemize}}
    
\end{document}